% Input from Hikita-rewrite.tex.
\section{A second proof of the dimension bound}\label{sec:bps}

% COMMENT: drafted 2026-10-14 from two unattributed notes Ben received,
% hikita_bps_only_proof.pdf and hall_coulomb_comparison.pdf (not in the
% repository).  Credit is still to be settled.  Citation checks against the
% sources are recorded in NOTES.md, "The BPS route".

This section gives a second proof of Proposition~\ref{prop:loopremoval}:
$\dim B_0=\dim H^*(X_{\nu_{\det}})$ for every quiver.  It adds no loops and
does not use \S\ref{sec:loops}.  Instead, it writes $B_0$ as a quotient of a
space $P_\Bd$: the part of dimension $(\bv,1)$ of a cohomological Hall algebra
of the framed quiver $\hat Q$, modulo decomposable elements and the image of
one tautological operator.  BPS cohomology spans $P_\Bd$, and Botta and
Davison identify it with $H^*(X_{\nu_{\det}})$.  The argument gives the
dimension only, not the Poincar\'e polynomial of Theorem~\ref{thm:loops}.

The proof is conditional.  It uses results of Jindal and Neguț
\cite{jindalLoopNilpotent2026} whose proofs are only sketched there; see
Remark~\ref{rem:bps-status}.  Throughout, $F$ is trivial, $\hbar=0$ on the
Coulomb side, $X=X_{\nu_{\det}}$, and $B_0=A_0/I_B$, where $I_B$ is generated
by the elements of nonzero charge.

\subsection{The framed quiver}\label{ss:bps-framed}
Let $\hat Q$ be $Q$ with one more vertex $\infty$ and $w_i$ arrows from
$\infty$ to $i$, as in \S\ref{ss:setup}, and put $\Bd=(\bv,1)$.  Let
$V_\infty=\C$ and $\hat G=G\times\GL(V_\infty)$.  After choosing bases of the
$W_i$, the representation space of $\hat Q$ of dimension $\Bd$ is $N$, with
$\Hom(W_i,V_i)=\Hom(V_\infty,V_i)^{w_i}$.  Let $\hat A$ be the Coulomb branch
of $(\hat G,N)$, at $\hbar=0$ and without flavor.  Its charges lie in
$\Z^{Q_0}\times\Z$.  The scalars in $\hat G$ act trivially on $N$, and this
splits $\hat A$, as in \pre{(8.4)}.

We use the abelianization of \S\ref{sec:deform}, as set out before
Proposition~\ref{prop:loopremoval}, for $\hat Q$ and $\Bd$ as well as for $Q$
and $\bv$.  Put $\hat\Omega=\Omega\sqcup\{\infty\}$, and let $\hat x_d$,
$d\in\hat\Omega$, be the coordinates on the Lie algebra of the diagonal torus
$\hat T$ of $\hat G$.  Put
\[
 x_d=\hat x_d-\hat x_\infty\quad(d\in\Omega),\qquad x_\infty=\hat x_\infty,
\]
so that $H^*_{\hat T}(\mathrm{pt})=\Pol[x_\infty]$.  Let $\hat\one\in\Z^{\hat\Omega}$
be the coweight with all entries $1$, and put
\[
 b=\sum_{d\in\Omega}x_d\in R,\qquad D=1+\sum_{k\in Q_0}v_k .
\]
Write $\hat r_\lambda$ and $\hat{\mathsf S}_\lambda(p)$ for the elements of
\S\ref{sec:deform} defined for $(\hat T,N)$, and $r_\lambda$, $\mathsf
S_\lambda(p)$ for those of $(T,N)$, written in the variables $x_d$.

\begin{lemma}\label{lem:bps-split}
Identify $e^{(\mu,0)}$ with $e^\mu$ for $\mu\in\Z^\Omega$.  Then:
\begin{enumerate}[leftmargin=2em]
\item $\hat r_{(\mu,k)}=r_{\mu-k\one}\,e^{k\hat\one}$ for $(\mu,k)\in\Z^\Omega\times\Z$;
in particular $\hat r_{\pm\hat\one}=e^{\pm\hat\one}$;
\item $\hat A=A_0[x_\infty,e^{\pm\hat\one}]$, and the map
$A_0\otimes\C[x_\infty,e^{\pm\hat\one}]\to\hat A$ is an isomorphism;
\item the part of $\hat A$ of charge $(\gamma,k)$ is
$A_0^{(\gamma-k\bv)}[x_\infty]\,e^{k\hat\one}$;
\item $\sum_{d\in\hat\Omega}\hat x_d=b+Dx_\infty$.
\end{enumerate}
\end{lemma}

\begin{proof}
The $\hat T$-weights of $N$ are the $\hat x_e-\hat x_d=x_e-x_d$, for the arrows
of $Q$, and the $\hat x_d-\hat x_\infty=x_d$, for the framing.  In the
variables $x_d$ these are the $T$-weights $\cB$ of $N$, and each pairs with
$(\mu,k)$ as it pairs with $\mu-k\one$.  This gives~(1), as
$e^{(\mu,k)}=e^{(\mu-k\one,0)}e^{k\hat\one}$.  Since $V_\infty$ is a line,
the root factors also agree: $\mathrm{eu}_{(\mu,k)}=\mathrm{eu}_{\mu-k\one}$.
So $\hat{\mathsf S}_{(\mu,k)}(p)=\mathsf S_{\mu-k\one}(p)\,e^{k\hat\one}$ for
$p\in\Pol^{\cW_\mu}[x_\infty]$.

A coweight $(\mu,k)$ has all entries in $\{0,1\}$ exactly when
$k\in\{0,1\}$ and $\mu-k\one$ has all entries in $\{0,1\}$ (if $k=0$) or in
$\{-1,0\}$ (if $k=1$); similarly for $\{0,-1\}$.  Also
$H^*_{\hat G}(\mathrm{pt})=R[x_\infty]$.  So~\ref{I:gen}, for $\hat Q$ and for $Q$,
shows that $\hat A$ is generated over $R[x_\infty]$ by the
$\mathsf S_\nu(p)\,e^{k\hat\one}$, with $\nu$ minuscule and $k\in\{0,\pm1\}$,
and that these generate $A_0[x_\infty,e^{\pm\hat\one}]$.  This is the first
part of~(2).  The map in~(2) is injective because
$\Z^{\hat\Omega}=\Z^\Omega\times\Z\hat\one$ and $x_\infty$ is transcendental
over $\mathrm{Frac}(\Pol)$.  For~(3), the coweight $(\mu,k)$ has charge
$(c,k)$, where $c$ is the charge of $\mu$, and $\mu-k\one$ has charge
$c-k\bv$.  For~(4), $\sum_{d\in\Omega}(x_d+x_\infty)+x_\infty=b+Dx_\infty$.
\end{proof}

Define an algebra homomorphism
\[
 \Theta\colon\hat A\to B_0,\qquad a\,x_\infty^je^{k\hat\one}\mapsto[a]\,(-b/D)^j\qquad(a\in A_0).
\]
It is well defined by Lemma~\ref{lem:bps-split}(2).  It is an algebra map
because the isomorphism there is one of algebras, and $\Theta$ is the tensor
product of the quotient map $A_0\to B_0$ with
$x_\infty\mapsto-b/D$, $e^{\hat\one}\mapsto1$.
Lemma~\ref{lem:bps-split}(3),(4) give:
\begin{enumerate}[leftmargin=2.6em,label=($\Theta$\arabic*)]
\item\label{Th:charge} $\Theta$ kills the part of $\hat A$ of charge
$(\gamma,k)$ unless $\gamma=k\bv$;
\item\label{Th:det} $\Theta\big((\sum_{d\in\hat\Omega}\hat x_d)\,a\big)=0$ for all
$a\in\hat A$;
\item\label{Th:coef} $\Theta\big(p(\hat x_d-\hat x_\infty:d\in\Omega)\big)=[p]$ for
$p\in R$.
\end{enumerate}
The value $-b/D$ is the only one compatible with~\ref{Th:det}.  The
corresponding projection $\hat x\mapsto\hat x-\frac1D(\sum_d\hat x_d)\hat\one$
splits off the scalars in the Lie algebra of $\hat T$.  Since it divides by
$D$, it does not come from a splitting of $\hat T$, which is why we use
$\Theta$ and not a decomposition of $\hat G$.

\begin{lemma}\label{lem:bps-charge}
If $0\neq\alpha\le\Bd$ and $\alpha\neq\Bd$, write $\alpha=(\gamma,k)$.  Then
$\gamma\neq k\bv$.
\end{lemma}

\begin{proof}
Here $k\in\{0,1\}$.  If $k=0$ then $\gamma\neq0$.  If $k=1$ then
$\gamma\le\bv$ and $\gamma\neq\bv$.
\end{proof}

\subsection{The loop-nilpotent Hall algebra}\label{ss:bps-hall}
Let $\mathsf P=\C[\hbar,u_a:a\in\hat Q_1]$, with maximal ideal
$\fm=(\hbar,u_a)$.  Here $\hbar$ is a new parameter on the Hall algebra side; it
becomes the loop rotation parameter of the Coulomb branch, which we then set
to $0$.  The tripled quiver of $\hat Q$ has the arrows $a$ of $\hat Q$, their
reverses $a^*$, and a loop $\omega_k$ at every vertex $k$.  It carries its
canonical potential and the action of a torus with weights $-u_a$ on $a$,
$u_a-\hbar$ on $a^*$ and $\hbar$ on $\omega_k$, which preserves the potential
\cite[\S3.1]{jindalLoopNilpotent2026}.  Let
$\cH_{\mathsf P}=\bigoplus_{\Bn}\cH_{\mathsf P,\Bn}$ be the loop-nilpotent
cohomological Hall algebra of \cite[\S3]{jindalLoopNilpotent2026}, with
coefficients extended to $\C$.  It is the equivariant critical cohomological
Hall algebra of the tripled quiver, with the $\omega_k$ required to be
nilpotent.  It is graded by dimension vectors
$\Bn\in\Z_{\ge0}^{\hat Q_0}$, with $\cH_{\mathsf P,0}=\mathsf P$, and each
$\cH_{\mathsf P,\Bn}$ is a module over
$H^*_{G_\Bn}(\mathrm{pt})\otimes\mathsf P=\mathsf P[z_{k,a}:a\le n_k]^{\mathrm{sym}}$, where
$G_\Bn=\prod_k\GL_{n_k}$.  We use the product $F\star G$, the product of
\cite{jindalLoopNilpotent2026} twisted by the sign
$(-1)^{\sum_{a\colon k\to l}n_lm_k}$ for $F$ of dimension $\Bn$ and $G$ of
dimension $\Bm$, as in \pre{(8.1)}.  The twist changes a word only by a sign
depending on the dimensions of its letters, so it does not change spans of
words.

For $g\in\mathsf P[z]^{\mathrm{sym}}$ let $e_\Bn(g)=g\cdot[\omega=0]_\Bn\in\cH_{\mathsf
P,\Bn}$ \cite[(98)]{jindalLoopNilpotent2026}.  Under the shuffle map of
\cite[\S3]{jindalLoopNilpotent2026} it goes to
$g(z)\prod_k\prod_{a,b=1}^{n_k}(z_{k,b}-z_{k,a}+\hbar)$
\cite[Lemma~3.6]{jindalLoopNilpotent2026}.  By definition,
$e_\Bn(hg)=h\cdot e_\Bn(g)$ for symmetric $h$.  Let $C_\Bn$ be the action of
$\sum_{k,a}z_{k,a}=c_1(\bigotimes_k\det V_k)$, the first Chern class of the
total determinant, on $\cH_{\mathsf P,\Bn}$.  It is a derivation:
\begin{equation}\label{eq:bps-derivation}
 C_{\Bn+\Bm}(F\star G)=(C_\Bn F)\star G+F\star(C_\Bm G),
\end{equation}
because the total determinant is multiplicative in short exact sequences;
\cite[Prop.~4.1]{davisonAffineBPS2022} proves this vertex by vertex.

Let $\cH=\cH_{\mathsf P}/\fm\cH_{\mathsf P}$, graded by the $\cH_\Bn$, and write
$e_\Bn(g)$ and $C_\Bn$ also for their images.  This specialization is
abstract.  The shuffle image of $e_\Bn(g)$ is divisible by $\hbar^{|\Bn|}$, so
it vanishes in the specialized shuffle algebra, but the class of $e_\Bn(g)$ in
$\cH$ need not vanish.  Put
\[
 \cT_\Bd=\sum_{\substack{\alpha+\beta=\Bd\\ \alpha,\beta\neq0}}\cH_\alpha\star\cH_\beta,
 \qquad
 P_\Bd=\cH_\Bd\big/(\cT_\Bd+C_\Bd\cH_\Bd).
\]
By~\eqref{eq:bps-derivation}, $C_\Bd\cT_\Bd\subseteq\cT_\Bd$.

We use the following inputs.
\begin{enumerate}[leftmargin=2.6em,label=(H\arabic*)]
\item\label{H:poly} \emph{Polynomial generation}
\cite[Thm.~2.4, Lemma~3.6, Prop.~3.7, Thm.~3.5]{jindalLoopNilpotent2026}.
$\cH_{\mathsf P}$ is generated as a $\mathsf P$-algebra by the elements
$e_\Bn(g)$.
\item\label{H:bps} \emph{BPS generation} \cite[(101)--(110),
Cor.~3.15]{jindalLoopNilpotent2026}.  There are $\mathsf P$-modules $L_\Bn$, the
loop-nilpotent BPS cohomology, and $\mathsf P$-linear maps
$\beta_\Bn\colon L_\Bn\to\cH_{\mathsf P,\Bn}$, such that $\cH_{\mathsf P}$ is
generated as a $\mathsf P$-algebra by the $C_\Bn^k\beta_\Bn(\ell)$, for $k\ge0$
and $\ell\in L_\Bn$.  This rests on cohomological integrality
\cite[Thm.~C]{davisonCohomologicalDonaldsonThomas2020}, in its equivariant
form \cite[Thm.~11.7]{davisonBPSAlgebras2023},
\cite[Thm.~3.4]{davisonAffineBPS2022}.  The tautological class $u$ of
\cite[(103)]{jindalLoopNilpotent2026} acts as $C_\Bn$; this is the
convention of \cite[\S3.4, (18)--(19)]{davisonAffineBPS2022}.
\item\label{H:framed} \emph{Framed BPS cohomology.}  If $\bw\neq0$, then
$L_\Bd$ is a free $\mathsf P$-module of rank $\dim H^*(X)$.  Indeed,
$L_\Bn\cong\hbar M_\Bn$ as $\mathsf P$-modules
\cite[Prop.~3.14]{jindalLoopNilpotent2026}, where $M_\Bn$ is the equivariant
BPS Lie algebra of the preprojective algebra of $\hat Q$ in dimension $\Bn$
\cite[Def.~3.6]{bottaOkounkovConjecture2023}.  Next,
$M_\Bn\cong M_\Bn^{\mathrm{noneq}}\otimes\mathsf P$
\cite[Cor.~11.10]{davisonBPSAlgebras2023}; the group allowed there contains
our torus.  Finally, $M_\Bd^{\mathrm{noneq}}\cong H^*(X)$, up to a shift of
degrees \cite[Cor.~3.13]{bottaOkounkovConjecture2023}.  The last result uses
framing arrows from $\infty$ to $i$, as here, and the stability $\delta_\infty$,
for which a module of dimension $\Bd$ is semistable exactly when no proper
submodule contains $V_\infty$.  So it is semistable exactly when the images of
the $W_i$ generate $V$, which is $\nu_{\det}$-stability by~\ref{I:cyc}.
\end{enumerate}

\begin{remark}[status of the inputs]\label{rem:bps-status}
With independent parameters, the genericity conditions of
\cite[\S2.2]{jindalLoopNilpotent2026} hold.
\begin{enumerate}[leftmargin=2em]
\item \cite[Thm.~2.4]{jindalLoopNilpotent2026} is a statement about a
shuffle algebra, and its proof is only sketched: it is ``completely
analogous'' to the K-theoretic case treated in earlier work of the same
authors.  The transfer to $\cH_{\mathsf P}$ uses
\cite[Prop.~3.7]{jindalLoopNilpotent2026} (proof sketched),
\cite[Lemma~3.6]{jindalLoopNilpotent2026} (stated without proof) and the
injectivity of the shuffle map \cite[Thm.~3.5]{jindalLoopNilpotent2026} (proof
omitted).
\item The PBW isomorphism for the loop-nilpotent algebra is
\cite[(109)]{jindalLoopNilpotent2026}, stated as an isomorphism of graded
vector spaces.  That it is induced by the product, which is what generation
needs, is implicit.  The proof of \cite[Cor.~3.15]{jindalLoopNilpotent2026}
cites the PBW theorem for the full algebra, Thm.~3.10 there, where (109) is
needed.
\item \cite[Prop.~3.14]{jindalLoopNilpotent2026} has a complete proof, which
uses \cite[Thm.~3.5]{jindalLoopNilpotent2026}.  The identification
\cite[Cor.~3.13]{bottaOkounkovConjecture2023} is sheaf-theoretic.  It combines
a comparison of BPS sheaves for different stability conditions
\cite[Prop.~3.4]{bottaOkounkovConjecture2023} with the description of the
stable framed BPS sheaf as the intersection complex of $X\times\A^1$
\cite[Prop.~3.11]{bottaOkounkovConjecture2023}, where $\A^1$ records the
common scalar of the loops $\omega_k$.  No Kac polynomials, point counts or
Hausel's formula enter.
\end{enumerate}
\end{remark}

\subsection{The Hall--Coulomb map}\label{ss:bps-map}
For $\Bn\le\Bd$ let $\lambda_\Bn\in\Z^{\hat\Omega}$ be the coweight that is $1$ on
the first $n_k$ elements of each $\hat\Omega_k$ and $0$ elsewhere.  It is
minuscule, with charge $\Bn$, and $\lambda_\Bd=\hat\one$.  For $g$ in
$\C[z_{k,a}:a\le n_k]^{\mathrm{sym}}$ write $g(\hat x_\Bn)$ for $g$
evaluated at the $\hat x_d$ with $(\lambda_\Bn)_d=1$.  It lies in
$\Pol[x_\infty]^{\cW_{\lambda_\Bn}}$.

\begin{lemma}[\pre{Lemma~8.1}]\label{lem:hall-coulomb}
Assume~\ref{H:poly}.  There is an algebra homomorphism
$\Upsilon\colon(\cH,\star)\to\hat A$ with the following properties.  It sends
$\cH_\Bn$ into the part of charge $\Bn$, and $e_\Bn(g)$ to $0$ if
$\Bn\not\le\Bd$.  For $\Bn\le\Bd$,
\[
 \Upsilon(e_\Bn(g))=\epsilon_\Bn\,\hat{\mathsf S}_{\lambda_\Bn}\big(g(\hat x_\Bn)\big),
\]
with $\epsilon_\Bn=\pm1$ depending only on $\Bn$.  In particular
$\Upsilon(e_\Bd(g))=\epsilon_\Bd\,g(\hat x)\,e^{\hat\one}$.
\end{lemma}

\begin{proof}[Sketch]
We follow \pre{Lemma~8.1}, which treats the shuffle algebra; it applies to
$\cH_{\mathsf P}$ through the shuffle map, an algebra homomorphism.  Over
$\mathrm{Frac}(\mathsf P)$, Tsymbaliuk's homomorphism from the shuffle
algebra to difference operators
\cite[Thm.~5.7]{tsymbaliukDifferenceOperators2023} gives an algebra
homomorphism $\Upsilon_{\mathsf P}$ from $(\cH_{\mathsf P},\star)$ to
difference operators in the $\hat x_d$, with
$D^\lambda f(\hat x)=f(\hat x+\hbar\lambda)D^\lambda$.  It sends $e_\Bn(g)$ to
$0$ if $\Bn\not\le\Bd$, and otherwise to $\pm E_\Bn(g)$, where
\[
 E_\Bn(g)=\sum_{w\in\cW/\cW_{\lambda_\Bn}}w\Big(g(\hat x_\Bn)\,
 \frac{\prod_{a\colon k\to l}\prod(\hat x_d-\hat x_e+u_a)}
      {\prod_k\prod(\hat x_d-\hat x_e)}\Big)D^{w\lambda_\Bn}
 \qquad\text{\pre{(8.2)}}.
\]
In the numerator $d\in\hat\Omega_k$ and $e\in\hat\Omega_l$; in the
denominator $d,e\in\hat\Omega_k$.  In both, $(\lambda_\Bn)_d=1$ and
$(\lambda_\Bn)_e=0$.  The reason is that Tsymbaliuk's map evaluates a shuffle
element on strings with step $\hbar$.  The factor $\prod(z_{k,b}-z_{k,a}+\hbar)$
vanishes on every string of length at least two, so only minuscule shifts
survive.  Its diagonal factors $\hbar^{|\Bn|}$ cancel the normalization of
\pre{Lemma~8.1}.

The coefficients of the $E_\Bn(g)$ lie in $\mathsf P[\hat x]$, localized at the
$\hat x_d-\hat x_e+r\hbar$ for $d\neq e$ in the same $\hat\Omega_k$ and
$r\in\Z$.  This ring is stable under the shifts, so the difference operators
with coefficients in it form a $\mathsf P$-algebra $\cD_{\mathsf P}$.  By~\ref{H:poly},
$\Upsilon_{\mathsf P}(\cH_{\mathsf P})\subseteq\cD_{\mathsf P}$, with no parameter
inverted.  Reducing modulo $\fm$ makes the $D^\lambda$ commute, and turns
$\cD_{\mathsf P}\otimes_{\mathsf P}\mathsf P/\fm$ into
$\Pol[x_\infty][(\hat x_d-\hat x_e)^{-1}][\Z^{\hat\Omega}]$, with
$D^\lambda\mapsto e^\lambda$.  Call the resulting map $\Upsilon$.

At $u_a=0$ the denominator of $E_\Bn(g)$ is $\mathrm{eu}_{\lambda_\Bn}$, and the
numerator is, up to a sign depending on $\Bn$, the product of the weights
$\beta$ of $N$ with $\langle\beta,\lambda_\Bn\rangle=-1$.  This is the
prefactor of $\hat r_{\lambda_\Bn}$.  So $\Upsilon(e_\Bn(g))$ is
$\pm\hat{\mathsf S}_{\lambda_\Bn}(g(\hat x_\Bn))$, which is the image of a dressed
minuscule monopole \cite[Prop.~6.6]{BFN}.  The images of the generators lie
in $\hat A$, and by~\ref{H:poly} so does the image of $\cH$.  The charge of
$w\lambda_\Bn$ is $\Bn$.  For $\Bn=\Bd$ both products are empty and
$\cW_{\lambda_\Bd}=\cW$, so $E_\Bd(g)=g(\hat x)D^{\hat\one}$.
\end{proof}

% COMMENT: [plausible, not checked against Tsymbaliuk's formula] Step 3
% below can avoid (H1).  At hbar = 0 the image of h F, for h symmetric,
% multiplies the D^lambda term of the image of F by h evaluated on the
% strings, so the image of C_d F is delta(image of F), where delta multiplies
% e^lambda by <lambda, \hat x>; that is, delta is minus the Poisson bracket
% with the Casimir (1/2) sum \hat x_d^2.  Then Theta(delta(a)) = 0 for a of
% charge d: delta(e^{\hat 1}) = (sum \hat x_d) e^{\hat 1}, and on A_0^{(0)} =
% R + sum A^{(-c)}A^{(c)} (Lemma lem:generic) delta kills R and preserves I_B.
% This is the derivation D_{i,k} of the preprint's section 5.  (H1) would
% then be needed only for the image of Upsilon to lie in \hat A.

\subsection{The bound}\label{ss:bps-bound}

\begin{theorem}\label{thm:bps}
Assume~\ref{H:poly}, \ref{H:bps} and~\ref{H:framed}.  For every quiver $Q$,
loops and multiple arrows allowed, and all $\bv$ and $\bw$,
\[
 \dim B_0\le\dim H^*(X).
\]
\end{theorem}

\begin{proof}
If $\bv=0$, then $B_0=\C$ and $X$ is a point.  If $\bw=0$ and $\bv\neq0$,
then $X$ is empty by~\ref{I:cyc}.  Also $\one$ pairs to zero with every weight
of $N$, so $r_{\pm\one}=e^{\pm\one}$, and by~\ref{I:gen} these lie in $A_0$.
They are inverse to each other and have nonzero charge $\pm\bv$, so $1\in I_B$
and $B_0=0$.  So assume $\bv\neq0$ and $\bw\neq0$, and let
\[
 \psi=\Theta\circ\Upsilon|_{\cH_\Bd}\colon\cH_\Bd\to B_0 .
\]

\emph{Step 1: $\psi$ is onto.}  For $p\in R$, let $g=p(z_{k,a}-z_{\infty,1})$,
which is symmetric at each vertex.  By Lemma~\ref{lem:hall-coulomb},
$\Upsilon(e_\Bd(g))=\epsilon_\Bd\,p(\hat x_d-\hat x_\infty)\,e^{\hat\one}$, so
$\psi(e_\Bd(g))=\epsilon_\Bd[p]$ by~\ref{Th:coef}.  The map $R\to B_0$ is onto
by Lemma~\ref{lem:generic}.

\emph{Step 2: $\psi(\cT_\Bd)=0$.}  Let $\alpha+\beta=\Bd$ with
$\alpha,\beta\neq0$.  Then $\Upsilon(\cH_\alpha)$ has charge $\alpha$, which
$\Theta$ kills by~\ref{Th:charge} and Lemma~\ref{lem:bps-charge}.  As $\Upsilon$
and $\Theta$ are multiplicative, $\psi(\cH_\alpha\star\cH_\beta)=0$.

\emph{Step 3: $\psi(C_\Bd\cH_\Bd)=0$.}  By~\ref{H:poly}, $\cH$ is spanned by
words in the $e_\Bn(g)$.  Letters with $\Bn=0$ are scalars, and a word of
dimension $\Bd$ with at least two letters of nonzero dimension lies in
$\cT_\Bd$.  So $\cH_\Bd$ is spanned by the $e_\Bd(g)$ and $\cT_\Bd$.  Now
$C_\Bd e_\Bd(g)=e_\Bd\big((\sum_{k,a}z_{k,a})g\big)$, so
\[
 \psi(C_\Bd e_\Bd(g))=\epsilon_\Bd\,\Theta\Big(\big(\textstyle\sum_{d\in\hat\Omega}\hat x_d\big)\,g(\hat x)\,e^{\hat\one}\Big)=0
\]
by~\ref{Th:det}.  Also $C_\Bd\cT_\Bd\subseteq\cT_\Bd$, which $\psi$ kills by
Step~2.  So $\psi$ induces a surjection $P_\Bd\to B_0$.

\emph{Step 4: the bound.}  Put $\bar L_\Bd=L_\Bd/\fm L_\Bd$.  By~\ref{H:bps},
and as tensor products are right exact, $\cH$ is generated by the images of
the $C_\Bn^k\beta_\Bn(\ell)$.  The word argument of Step~3 shows that $\cH_\Bd$
is spanned by $\cT_\Bd$ and the $C_\Bd^k\beta_\Bd(\ell)$, $\ell\in\bar L_\Bd$.
The terms with $k\ge1$ lie in $C_\Bd\cH_\Bd$, so $\bar L_\Bd\to P_\Bd$ is onto.
By~\ref{H:framed}, $\dim\bar L_\Bd=\dim H^*(X)$.  Hence
\[
 \dim B_0\le\dim P_\Bd\le\dim\bar L_\Bd=\dim H^*(X).\qedhere
\]
\end{proof}

No injectivity is needed anywhere in this proof, and no map is compared with
a geometric one.

\begin{corollary}\label{cor:bps}
Assume~\ref{H:poly}, \ref{H:bps} and~\ref{H:framed}.  Then
Proposition~\ref{prop:loopremoval} holds: for every quiver,
$\Phi\colon B_0\to H^*(X)$ is an isomorphism.  For $\bw\neq0$ the
surjections $\bar L_\Bd\to P_\Bd\to B_0$ of the proof are isomorphisms of
vector spaces.  With this in place of Proposition~\ref{prop:loopremoval}, the
proof of Theorem~\ref{thm:classical}, and so of Theorem~\ref{thm:main}, goes
through without \S\ref{sec:loops}.
\end{corollary}

\begin{proof}
The surjection $\Phi\colon B_0=B^{\nu_{\det}}_0\to H^*(X)$ of \S\ref{sec:map}
and Theorem~\ref{thm:bps} give $\dim B_0=\dim H^*(X)$, finite, so $\Phi$ is an
isomorphism.  Then every map in the chain $\bar L_\Bd\to P_\Bd\to B_0$ is a
surjection between spaces of the same finite dimension.
\end{proof}

\begin{remark}\label{rem:bps-compare}
Compared with the proof through \S\ref{sec:loops}:
\begin{enumerate}[leftmargin=2em]
\item This proof uses \S\ref{sec:map}, Lemmas~\ref{lem:flavor}
and~\ref{lem:betti}, and inputs~\ref{I:gen}--\ref{I:cyc}.  It does not use
\S\ref{sec:loops}, the loops added in Proposition~\ref{prop:loopremoval}, or
inputs~\ref{I:hypertoric}--\ref{I:core}.
\item It needs~\ref{H:poly}--\ref{H:framed} and \pre{Lemma~8.1}.  The proof in
the rest of the paper does not use Jindal--Neguț at all, because
Lemma~\ref{lem:generic} replaces \pre{\S8}.
\item It gives $\dim H^*(X)$, but not the formula of Theorem~\ref{thm:loops} or
the multipartition filtration.  It does not equip $\bar L_\Bd$ or $P_\Bd$ with
the ring structure of $B_0$, and we have not checked whether the maps
$\bar L_\Bd\to P_\Bd\to B_0$ respect gradings.
\end{enumerate}
\end{remark}
